\section{Results and Discussion}

\subsection{Forecasting Insights}

The demand analysis of the forecast shows that there are a few crucial trends in platform-level demand. Among the main findings is the existence of apparent differences in the weekly demand, as the number of demand differs by the day of the week. An example is that some days record higher order volumes and others have low demand. This is an indication of customer behavior and habitual consumption. 

Furthermore, the data has seasonal trends. The forecasting model captures the repeating patterns like weekly seasonality, meaning that demand has predictable patterns over time. The model can recognize the trends within the historical data and give consistent predictions, even though there are short-term changes and volatility. 

These observations prove that demand is not haphazard but capable of following specific patterns, and these can be used to plan and make decisions in a more efficient way.


\subsection{Optimization Results}

The recommendations for the optimization model are related to the allocation of orders and promotions among delivery platforms so as to maximize the profit. The findings show that each platform has a varying impact on overall business performance, as a result of variations in customer demand, platform commissions, and sales volume.

The optimization framework helps identify allocation strategies that focus on the more rewarding platforms while taking into account the expected demand constraints and operational limitations. The optimized allocation strategy is more effective in using business resources and marketing efforts than the equal allocation or the intuition-based allocation strategies.

The implementation of the optimization model based on the results of simulation and back testing, may lead to the potential for the improvement of profitability from historical decision making processes. The forecasting outputs are incorporated into the optimization process, which can help to minimize uncertainty caused by demand fluctuations and make more informed managerial decisions.

The results also show that forecast-based optimization can augment the promotional planning process by optimizing campaigns to the periods and platforms where more is expected to be gained from the promotion. This helps businesses to make the best use of their marketing budgets and boost their overall efficiency.

Yet another lesson to be learned is that the optimization model enables strategic platform assessment. The framework allows companies to assess the performance of the platforms apart from the number of sales to also measure the profitability contribution, thus further adding to the business performance. This analysis can help management determine which platforms are more valuable to the business and where investments and promotions should be made.

The optimization results indicate that business efficiency can be increased, strategic planning can be supported and operational performance can be optimized in a third-party delivery environment by combining forecasting with profit-oriented decision models.


\subsection{Business Implications}

The suggested decision support structure offers a number of implications for HungryHub Oy. The company can shift from reacting to making decisions towards a more proactive and data-driven approach to operations, by combining data analysis, forecasting, and optimization into a single analytical model.

This framework will enable the company to have a better understanding of the behaviour of customer demand; platform specific performance, and future sales. These insights help to make better plans for promotion, pricing and resource planning decisions for various delivery platforms.

Besides, the forecasting aspect helps management to forecast the change in demand and accordingly prepare operational resources. This could help in optimizing inventory planning, manpower utilization, and campaign scheduling, especially during peak seasons. Better visibility of forecasting could also help to mitigate operational uncertainty and avoid excess costs of over-preparing or underutilizing resources.

The optimization component helps improve profits through the evidence-informed allocation decisions. This framework allows for decisions to be made based on forecasted demand patterns and for a profitability analysis of the platform. This is not merely relying on managerial intuition. This helps to use marketing resources and operate more efficiently.

The framework is also useful in the long term business planning process. The company can constantly track the performance of the platform and sales trends, which can help them understand the preferences of their users and see how well their platform partnerships are working. This can help management to make adjustments to business strategies in the fast-expanding third-party delivery market.

The overall system showcases the potential of DDD to enhance operational efficiency, minimize uncertainty, and boost competitiveness in the digital food delivery landscape.

\subsection{Key Findings}
In this study, the main findings from the forecasting, optimization, and business analysis are summarized in Table \ref{tab:key_findings}. The results reveal key patterns of demand, platform level performance differences, and the opportunities for using forecasting and optimization techniques as part of business decisions. The summary also illustrates how the proposed framework can help enhance operational efficiency, improve profitability and support strategic planning in third party delivery settings.
\begin{table}[h]
\centering
\footnotesize
\renewcommand{\arraystretch}{1.6}
\setlength{\tabcolsep}{4pt}
\begin{tabular}{|p{4cm}|p{5.5cm}|p{4.5cm}|}
\hline
\textbf{Key Finding} & \textbf{Observation} & \textbf{Business Implication} \\ \hline

Weekly demand variation &
Customer demand differs across weekdays and weekends &
Supports better staffing and operational planning \\ \hline

Seasonal demand patterns &
Repeating sales trends were identified in historical data &
Enables improved forecasting and inventory preparation \\ \hline

Platform-level performance differences &
Wolt and Foodora showed different sales and profitability behaviors &
Assists in platform-specific strategic planning \\ \hline

Forecasting model effectiveness &
The forecasting model captured major demand trends with acceptable accuracy &
Supports data-driven demand prediction and planning \\ \hline

Optimization-based allocation &
Profit-oriented allocation produced better outcomes than historical allocation strategies &
Improves resource utilization and profitability \\ \hline

Forecast-driven decision-making &
Combining forecasting with optimization improved managerial insight &
Reduces uncertainty in operational decisions \\ \hline

Promotional planning opportunities &
High-demand periods can be identified in advance &
Enables more efficient promotional campaign scheduling \\ \hline

Data integration benefits &
Centralized analysis improved visibility across sales channels &
Supports comprehensive business monitoring \\ \hline

Strategic business value &
The framework supports proactive rather than reactive decision-making &
Enhances competitiveness in the delivery market \\ \hline

Scalability potential &
The framework can be extended into a real-time decision-support system &
Provides opportunities for future system development \\ \hline

\end{tabular}
\caption{Key Findings and Business Implications}
\label{tab:key_findings}
\end{table}

\subsection{Research Questions and Answers}

This section presents the main research questions of the study and explains how each question has been addressed based on the results.

\textbf{RQ1: How can sales data from third-party delivery platforms be structured into a unified analytical dataset?}

The study addressed this question by cleaning and integrating data from Wolt and Foodora into a single dataset. This unified dataset enabled consistent analysis across platforms and served as the foundation for forecasting and optimization.

\textbf{RQ2: Which forecasting approach provides the most accurate predictions of platform-level demand?}

The Prophet model was selected and applied due to its ability to capture trend and seasonal patterns effectively. The evaluation results showed that Prophet provides reliable and stable forecasts suitable for business decision-making.

\textbf{RQ3: How can demand forecasts be transformed into profit-oriented decision recommendations?}

The study developed an optimization model that uses forecasted demand as input. This model generates recommendations on how to allocate orders or promotional efforts across platforms to maximize profit under given constraints.

\textbf{RQ4: What is the potential business impact of a prescriptive decision-support framework in optimizing platform sales performance?}

The simulation and back-testing results indicate that the proposed framework can improve decision-making and increase potential profitability. By combining forecasting and optimization, the system provides more effective and data-driven recommendations compared to traditional approaches.